\documentclass[../main.tex]{subfiles}
\begin{document}

\section{Lý do chọn đề tài}
Các nhóm phát triển phần mềm nhỏ---đặc biệt là nhóm sinh viên---thường quản lý công việc bằng tin
nhắn, bảng tính hoặc một công cụ đa năng như Trello hay Jira. Mỗi lựa chọn đều có cái giá của nó.
Tin nhắn làm thất lạc thông tin; bảng tính không biết gì về trạng thái, phụ thuộc hay hạn chót; còn
các công cụ lớn thì nặng, phụ thuộc chuột, luôn cần kết nối mạng và thu phí đúng những chức năng nhóm
cần như báo cáo Scrum. Trong khi đó, hai xu hướng giúp một công cụ tốt hơn trở nên khả thi: trình
duyệt giờ có thể lưu toàn bộ bản sao làm việc của dự án ngay trên máy (ý tưởng \emph{local-first}
\cite{kleppmann2019localfirst}), và các mô hình ngôn ngữ lớn có thể biến văn bản tự do---biên bản
họp, một mục tiêu còn mơ hồ---thành dữ liệu có cấu trúc, có thể hành động được.

Đề tài xây dựng một hệ thống quản lý dự án nhóm quanh ba ý tưởng: mọi thao tác đều làm
được bằng bàn phím; bảng công việc vẫn hoạt động khi mất mạng; và AI đảm nhận các phần tẻ nhạt của
quản lý (tóm tắt, biên bản, giao việc, chia nhỏ công việc, retrospective cuối sprint) nhưng luôn
trả về kết quả có cấu trúc, kiểm tra được.

\section{Mục tiêu}
\begin{itemize}
  \item Xây dựng ứng dụng web cho phép người dùng đăng ký, tạo dự án và mời thành viên, với vai trò
        (Quản lý dự án hoặc Thành viên) được gán \emph{theo từng dự án} chứ không theo tài khoản.
  \item Cung cấp bảng Kanban 4 cột và chế độ xem dạng bảng, điều khiển hoàn toàn bằng bàn phím, có
        phụ thuộc giữa các task và đường găng được tính tự động.
  \item Hỗ trợ quy trình Scrum: lập kế hoạch sprint, product backlog, bắt đầu và kết thúc sprint,
        biểu đồ burndown và velocity, retrospective do AI viết.
  \item Tích hợp mô hình ngôn ngữ lớn (Groq, model \code{openai/gpt-oss-120b}) cho năm chức năng AI,
        với cơ chế dự phòng ba tầng để ứng dụng không bao giờ lỗi khi dịch vụ AI không sẵn sàng.
  \item Giữ dự án cá nhân dùng được đầy đủ khi ngoại tuyến, đồng thời ngăn chỉnh sửa xung đột khi
        ngoại tuyến ở dự án có nhiều thành viên.
  \item Triển khai bản chạy cấu hình production trên tên miền công khai để nhóm kiểm thử, và kiểm
        chứng hệ thống bằng kiểm thử tự động.
\end{itemize}

\section{Phạm vi}
Hệ thống hướng tới nhóm nhỏ (khoảng 2--15 người) làm dự án phần mềm. Đây là ứng dụng web dùng được
trên trình duyệt máy tính và điện thoại, bằng tiếng Anh và tiếng Việt. Ngoài phạm vi: chấm công,
thanh toán, biểu đồ Gantt, ứng dụng di động gốc và tích hợp với các hệ thống quản lý issue khác.

\section{Phương pháp thực hiện}
Nhóm áp dụng chính quy trình mà sản phẩm hỗ trợ. Công việc được tổ chức theo sprint với một backlog
chung; mỗi thay đổi đi kèm kiểm thử, mọi quyết định quan trọng đều được ghi lại. Dự án duy trì tám
tài liệu sống (D1 Yêu cầu, D2 Bản đồ module, D3 Kiến trúc, D4 Hợp đồng API và dữ liệu, D5 Kiểm thử và
tiêu chí chấp nhận, D6 Rủi ro và chính sách, D7 Sổ tay phát triển và nhật ký quyết định, D8 Ma trận
truy vết yêu cầu) nằm cùng kho mã nguồn, để mỗi thay đổi mã và tài liệu mô tả nó được xem xét cùng
nhau. Kiểm thử theo hướng \emph{characterisation-first}: trước khi sửa đoạn mã chưa có test, viết
test mô tả hành vi hiện tại; bản thân các test được kiểm tra bằng \emph{kiểm thử đột biến} (mutation
testing)~\cite{demillo1978mutation}---cố ý làm hỏng mã và xác nhận có test báo lỗi.

\section{Tổ chức nhóm và phân công nhiệm vụ}
Nhóm gồm ba thành viên lớp Kỹ thuật máy tính K23A. Bảng~\ref{tab:tasks} trình bày phân công công việc.

\begin{longtable}{|p{3.2cm}|p{2.4cm}|p{8.0cm}|}
\caption{Phân công nhiệm vụ các thành viên}\label{tab:tasks}\\
\hline
\textbf{Thành viên} & \textbf{Vai trò} & \textbf{Nhiệm vụ} \\
\hline
\endfirsthead
\hline
\textbf{Thành viên} & \textbf{Vai trò} & \textbf{Nhiệm vụ} \\
\hline
\endhead
Phạm Văn Huynh \newline (KTMT K23A) & Nhóm trưởng \newline Scrum Master &
\begin{itemize}[leftmargin=*,topsep=0pt]
 \item Lập kế hoạch dự án, tổ chức sprint, theo dõi tiến độ
 \item Phân tích yêu cầu và bộ tài liệu D1--D8
 \item Module Scrum: vòng đời sprint, lập kế hoạch backlog, burndown, velocity, retrospective AI
 \item Tích hợp AI với Groq, thiết kế prompt; AI trả lời theo ngôn ngữ giao diện
 \item Triển khai: Docker, MySQL, Cloudflare Tunnel trên \code{koshi.huynhca.com}
 \item Dữ liệu kiểm thử, smoke test đầu-cuối, video demo, báo cáo
\end{itemize} \\
\hline
Phạm Minh Tú \newline (KTMT K23A) & Lập trình viên chính \newline (kiến trúc, frontend) &
\begin{itemize}[leftmargin=*,topsep=0pt]
 \item Kiến trúc tổng thể và lựa chọn công nghệ
 \item Frontend Vue 3: Pinia store, Kanban và bảng, bộ xử lý phím tắt
 \item Đồ thị phụ thuộc, sắp xếp topo và đường găng
 \item Lưu trữ local-first (IndexedDB) và chính sách ghi khi ngoại tuyến
 \item Trang giới thiệu, trang hồ sơ, đa ngôn ngữ Anh/Việt
\end{itemize} \\
\hline
Đàm Đức Đôn \newline (KTMT K23A) & Lập trình viên \newline (backend, chất lượng) &
\begin{itemize}[leftmargin=*,topsep=0pt]
 \item Backend FastAPI: xác thực, dự án, vai trò theo dự án và lời mời
 \item Bình luận, @mention, trả lời, đính kèm, thông báo, lịch sử task
 \item Lược đồ CSDL, Alembic migration, chuyển từ SQLite sang MySQL
 \item Tích hợp liên tục (GitHub Actions), bộ test backend và frontend
 \item Logic dự phòng AI và xoay vòng API key
\end{itemize} \\
\hline
\end{longtable}

\section{Cấu trúc báo cáo}
Ngoài phần mở đầu và kết luận, báo cáo gồm bốn chương. Chương~1 trình bày cơ sở lý thuyết: Kanban,
Scrum, đồ thị phụ thuộc và đường găng, phần mềm local-first, mô hình ngôn ngữ lớn và các công nghệ sử
dụng. Chương~2 phân tích yêu cầu. Chương~3 mô tả thiết kế hệ thống: kiến trúc, mô hình dữ liệu, API,
chuỗi dự phòng AI, bảo mật và chính sách ngoại tuyến. Chương~4 trình bày kết quả cài đặt kèm ảnh chụp
màn hình, chiến lược kiểm thử, kết quả và việc triển khai.

\end{document}
